% !TEX root = ../main.tex
\section{System Integration}

{\emergencystretch=3em \code{execute\_\allowbreak schedule(grid, charging\_\allowbreak station, schedule, algorithm)} connects the two parts.
\code{plan\_\allowbreak and\_\allowbreak execute} first runs \code{backtracking\_search} on the parsed constraints. It then
converts the assignment into the \code{schedule}: the task labels sorted by slot, so empty slots are
skipped. Each label is mapped to its cell through \code{grid.dirty}. The integration only uses the
outputs of the CSP and of \code{search}; the visiting order and the paths are never hard-coded.\par}

\begin{algorithm}[!htbp]
\caption{Executing the CSP schedule with graph search}
\begin{algorithmic}[1]
\Function{Execute-Schedule}{$grid, charger, schedule, algorithm$}
  \State $stops \gets [\,grid.dirty[\ell] \text{ for } \ell \in schedule\,] + [charger]$
  \State $current \gets charger$; \ $complete\_path \gets [charger]$; \ $total \gets 0$
  \For{$target \in stops$}
    \State $(path, cost, nodes) \gets \Call{Search}{grid, current, target, algorithm}$
    \If{$path = \textit{None}$}
      \State \Return $(complete\_path, total, \textit{False})$ \Comment{failure, \dec{15}}
    \EndIf
    \State $complete\_path \gets complete\_path + path[1{:}]$ \Comment{drop repeated first cell}
    \State $total \gets total + cost$; \ $current \gets target$
    \If{$target \neq charger$} mark $target$ cleaned
    \EndIf
  \EndFor
  \State \Return $(complete\_path, total, \textit{True})$ \Comment{all cleaned, back at charger}
\EndFunction
\end{algorithmic}
\end{algorithm}

The grid is never changed; cleaned cells are tracked separately (\dec{06}). This way UCS and
\astar{} can run one after the other on the same \code{Grid}. For each algorithm, the notebook output
shows the schedule, the visit sequence, each leg's path with its cost and nodes expanded, the final
grid (\code{*} path, \code{dk} cleaned), the total moves and cost, and the final status.

\textbf{Failure handling (\dec{15}).} If a leg has no path, execution stops at that leg. It returns
the path walked so far, its cost, and \code{success = False}. The status line names the unreachable
leg, the tasks cleaned and not cleaned, and states that the robot did not return to the charger, so
the system never claims success. If the CSP finds no schedule, no search is run: the path is
\code{[charger]}, the cost is 0, and the run is reported as a failure. Input errors are caught and
reported as \code{InputError} messages, so the notebook does not crash.
